
% !TEX TS-program = pdflatex
\documentclass[11pt]{article}

\usepackage[margin=1in]{geometry}
\usepackage{amsmath,amssymb,amsthm,mathtools}
\usepackage{enumitem}
\usepackage{hyperref}
\usepackage{microtype}

\hypersetup{
  colorlinks=true,
  linkcolor=blue,
  citecolor=blue,
  urlcolor=blue
}

\title{A Conceptual Primer on Giry--Kleisli de Finetti,\\ Measurability Bottlenecks, and Markov/Fortini Bridges\\(2026--02--27)}
\author{(AIrxiv note; conceptual summary for formalization work)}
\date{}

\newtheorem{definition}{Definition}
\newtheorem{theorem}{Theorem}
\newtheorem{lemma}{Lemma}
\newtheorem{remark}{Remark}

\newcommand{\NN}{\mathbb{N}}
\newcommand{\RR}{\mathbb{R}}
\newcommand{\ENN}{\mathbb{R}_{\ge 0}^{\infty}}
\newcommand{\PMeas}{\mathsf{P}}
\newcommand{\FinMeas}{\mathsf{M}_{<\infty}}
\newcommand{\ProbMeas}{\mathsf{PMeas}}
\newcommand{\Kern}{\mathsf{Kern}}
\newcommand{\Giry}{\mathsf{Giry}}
\newcommand{\id}{\mathrm{id}}
\newcommand{\Law}{\mathrm{Law}}
\newcommand{\E}{\mathbb{E}}

\begin{document}
\maketitle

\begin{abstract}
This note explains, at a strictly conceptual level, a formalization pipeline for
de Finetti-type representation theorems in the \emph{Giry monad} and its Kleisli category,
with a focus on a recurring bottleneck: upgrading \emph{pointwise} constructions of mediating
measure-valued maps into \emph{measurable kernels}. We also summarize how these ideas connect
to universal prediction perspectives (``Solomonoff-style'') and to Markov de Finetti
(Fortini/Diaconis--Freedman--type) proofs, highlighting the single crux that must be isolated:
\emph{cross-anchor conditional independence} of row-successor processes.
\end{abstract}

\tableofcontents

\section{Orientation: what is being proved?}

There are two parallel stories:

\begin{enumerate}[label=(\alph*)]
\item \textbf{Giry/Kleisli de Finetti.}
One models probabilistic programs as morphisms in the Kleisli category of the Giry monad:
objects are measurable spaces, and morphisms are \emph{Markov kernels} (measurable maps into
spaces of probability measures). The goal is to express ``exchangeable sources'' as mixtures
over latent parameters via a universal property: a \emph{limit cone} (or equivalently, a
universal mediator/factorization statement) in a suitable Kleisli setting.

\item \textbf{Markov/Fortini de Finetti for finite-state paths.}
One starts with a prefix law $\mu$ on finite words over $\mathrm{Fin}\;k$,
assumes Markov exchangeability (probability depends only on start state + transition counts),
extends $\mu$ to a path measure $P$ on $(\NN \to \mathrm{Fin}\;k)$, and aims to represent
cylinder probabilities as a mixture over Markov parameters (initial distribution + transition matrix),
\[
\mu(x_0\ldots x_n)=\int \mathrm{wordProb}(\theta, x_0\ldots x_n)\,d\pi(\theta).
\]
\end{enumerate}

In both stories, one wants to construct a \emph{latent object} (a parameter space, or a space of
directing measures/kernels) and show that all observed laws factor uniquely through it.

\section{Background: exchangeability and de Finetti in kernel language}

\begin{definition}[Exchangeability (informal)]
A stochastic process $(X_n)_{n\in\NN}$ is \emph{exchangeable} if its joint law is invariant under
finite permutations of indices. For Markov exchangeability (finite-state), invariance is weaker:
probabilities depend only on initial state and transition counts.
\end{definition}

Classically, de Finetti's theorem states: an exchangeable sequence is a mixture of i.i.d.\ sequences.
In \emph{kernel form}, ``mixture'' is a pushforward/integration over a latent parameter $\Theta$:
given a \emph{directing measure} $\pi$ on $\Theta$ and a kernel $K : \Theta \leadsto \Omega$,
the observed law is $\int K(\theta,\cdot)\,d\pi(\theta)$.

This suggests a categorical formulation: observed sources are induced from a universal latent object.
The formalization challenge is not the algebra of integrals---it is the measurability needed for the
kernels that implement the factorization.

\section{The Giry monad and its Kleisli category}

Let $\mathcal{M}$ be a category of measurable spaces and measurable maps. The Giry monad $\Giry$
maps a space $X$ to the space $\ProbMeas(X)$ of probability measures on $X$, equipped with the
\emph{evaluation $\sigma$-algebra} (the smallest $\sigma$-algebra making $\mu\mapsto \mu(A)$ measurable
for all measurable $A\subseteq X$).

\begin{definition}[Kleisli morphisms]
A Kleisli morphism $X \to_{\mathrm{Kl}(\Giry)} Y$ is a measurable map
$f : X \to \ProbMeas(Y)$.
Equivalently, this is a Markov kernel $f(x,\cdot)$ from $X$ to $Y$.
\end{definition}

Composition is given by integration (bind):
$(g \star f)(x, B) \;=\; \int g(y,B)\,df(x)(y)$.
The \emph{unit} sends $x$ to the Dirac measure $\delta_x$.

\section{Categorical de Finetti: the universal mediator picture}

A recurring template is:

\begin{enumerate}[label=\arabic*.]
\item Define an object $\Theta$ of latent parameters.
\item Define a kernel $K : \Theta \leadsto \Omega$ that generates i.i.d.\ sequences given $\theta$.
\item For an exchangeable source $P$ on $\Omega^\NN$, construct a directing measure $\pi$ on $\Theta$.
\item Prove the \emph{representation} $P = \int K(\theta,\cdot)\,d\pi(\theta)$ and a universal property:
any other latent representation factors uniquely through $\Theta$.
\end{enumerate}

In categorical terms, one expresses a cone from $P$ to a diagram of finite-marginal objects, and shows
it factors uniquely through a universal cone induced by $\Theta$.

\subsection{Why kernels (not just functions to measures) matter}

A frequent trap is:

\begin{quote}
``For each $y$, define a measure on $\Theta$ pointwise by some formula, so we have a mediator
$y \mapsto \pi_y$.'' \\
\emph{But} to be a Kleisli arrow, $y \mapsto \pi_y$ must be measurable into $\ProbMeas(\Theta).
$
\end{quote}

This is exactly where many proof attempts get stuck: \emph{measurability upgrades} from pointwise
definitions to genuine kernels.

\section{The measurability bottleneck: what usually goes wrong}

\subsection{Two $\sigma$-algebras on spaces of measures}

There are (at least) two natural measurable structures on spaces of measures:

\begin{enumerate}[label=(\alph*)]
\item The \textbf{Giry/evaluation} $\sigma$-algebra (generated by $\mu\mapsto \mu(A)$).
\item The \textbf{Borel} $\sigma$-algebra of a topology on measures (weak, weak-$*$, Prokhorov, \dots).
\end{enumerate}

When a proof uses topological continuity (e.g.\ continuity of a ``directing measure map'' in a metric
like L\'evy--Prokhorov), it naturally yields Borel measurability. But the Giry monad uses the
evaluation $\sigma$-algebra. To move between these worlds, one needs an inclusion/equality:
\[
\mathsf{borel}(\ProbMeas(X)) \;\le\; \mathsf{instMeasurableSpace}(\ProbMeas(X))
\quad\text{or}\quad
\mathsf{instMeasurableSpace}(\ProbMeas(X)) \;\le\; \mathsf{borel}(\ProbMeas(X)),
\]
depending on how each side is defined in the formalization.

\subsection{Why this is nontrivial in formalization}

Even if mathematicians ``know'' the relevant equalities under standard hypotheses
(Polish/standard Borel/Prokhorov theorems), the library may not contain the lemma in the needed form,
or it may require extra topological assumptions that are not currently available on your parameter
space.

\begin{remark}[A reliable separation principle]
It helps to isolate the exact missing implication as a \emph{single hypothesis} (e.g.\
$\mathsf{borel}\le \mathsf{instMeasurable}$), prove everything downstream from it, and then work
locally on proving the hypothesis under minimal standard conditions.
\end{remark}

\subsection{When measurable selection/disintegration is genuinely needed}

A frequent pattern is:

\begin{quote}
``We have existence and uniqueness of a conditional/directing object pointwise; therefore we can
choose it and proceed.''
\end{quote}

But \emph{classical choice} does not preserve measurability. If you need $y\mapsto \pi_y$ measurable,
you typically need either:
\begin{enumerate}[label=(\roman*)]
\item a measurable selection theorem (for set-valued maps),
\item a disintegration theorem (regular conditional probabilities),
\item or a proof that your pointwise definition is generated by countably many measurable coordinates.
\end{enumerate}

In practice, the weakest standard assumptions are those guaranteeing \emph{standard Borel} structure
on the spaces involved and existence of regular conditional probabilities.

\section{A workable design pattern: generated $\sigma$-algebras for parameter spaces}

One robust approach is to define the measurable structure on a parameter type $\Theta$ as the
\emph{least} $\sigma$-algebra that makes the evaluation maps you actually integrate against
measurable.

For a Markov parameter $\theta$ (initial distribution + transitions), one can generate a
$\sigma$-algebra by the \emph{word probability} maps
$\theta \mapsto \mathrm{wordProb}(\theta, xs)$ for all finite words $xs$.
This yields two benefits:

\begin{enumerate}[label=(\alph*)]
\item Every integrand $\mathrm{wordProb}(\theta, xs)$ is measurable \emph{by construction}.
\item To prove a map $g : \Omega^\NN \to \Theta$ is (AE-)measurable, it suffices to prove that all
compositions $\mathrm{wordProb}(\,g(\cdot)\,, xs)$ are (AE-)measurable.
\end{enumerate}

This is a powerful ``measurability-by-generators'' strategy: instead of fighting the library's
topological instances, you align the measurable structure with your proof needs.

\section{Markov/Fortini de Finetti: where the crux really lives}

For finite-state paths, a clean formalization often splits into:

\begin{enumerate}[label=\arabic*.]
\item \textbf{Row processes.} For each anchor state $i$, define a row-successor process that records
the successor state at the $n$th visit to $i$.
\item \textbf{Per-row de Finetti.} Show each row process is exchangeable (under suitable invariance),
therefore a mixture of i.i.d.\ sequences, giving a \emph{row kernel} $K_i(r)$.
\item \textbf{Lift to a Markov parameter.} Combine row kernels (and an initial kernel, often Dirac at
the start state) into a Markov parameter $\theta(\omega)$.
\item \textbf{Reconstruction.} Prove cylinder probabilities equal the integral of $\mathrm{wordProb}$
under the pushforward law of $\theta(\omega)$.
\end{enumerate}

Steps 1--3 are ``local'' and largely measurable-algebraic. Step 4 contains the true difficulty.

\subsection{The crux: cross-anchor conditional independence}

Per-row conditional i.i.d.\ gives, for each fixed $i$, a decomposition of the \emph{row marginal} law.
But reconstructing a full path cylinder typically needs a \emph{joint} statement tying together the
different anchors $i$ appearing in a word.

Conceptually, you need something like:

\begin{quote}
Conditional on the directing objects for each row, the row processes for different anchors behave
independently in the way required to factor a word probability into a product of row contributions.
\end{quote}

This is \emph{not} implied by ``each row is conditionally i.i.d.'' as a separate statement for each
$i$. One must either:
\begin{enumerate}[label=(\alph*)]
\item derive a joint law from stronger invariance/exchangeability hypotheses,
\item or assume an explicit joint law hypothesis (and label the theorem accordingly),
\item or redesign the proof to avoid needing a full joint independence statement (rare, but possible in
some finite-state combinatorial approaches).
\end{enumerate}

\subsection{A conceptual ``identity to aim for''}

For a word $a::b::xs$, a typical target identity is:
\[
P(\mathrm{cyl}(a::b::xs)) \;=\; \int \mathbf{1}\{ \omega_0=a\}\cdot
\prod_{t\ \text{transitions in }a::b::xs} K_{\text{row}(t)}(\omega)\big(\{\text{next}(t)\}\big)\; dP(\omega),
\]
i.e.\ cylinder probability equals the integral of a product of row-kernel singleton evaluations.
This is exactly the kind of statement often packaged as a \emph{cross-anchor product identity}.

\section{How the Giry story helps the Fortini story (and vice versa)}

The Giry/Kleisli development suggests two guiding principles that directly help the Fortini bridge:

\begin{enumerate}[label=\arabic*.]
\item \textbf{Isolate measurability from probability identities.}
Make the map $\omega \mapsto \theta(\omega)$ AE-measurable first (using generator-based
measurable spaces where possible), then reduce reconstruction to a purely measure-theoretic identity.

\item \textbf{Name the missing hypothesis honestly.}
If the formal library does not provide a theorem turning ``per-row conditional i.i.d.'' into the
needed cross-anchor product identity, the proof is stuck for a real mathematical reason, not an
automation failure. The remedy is to either strengthen assumptions or prove the missing theorem from
existing invariances (e.g.\ from Markov exchangeability of the original prefix law, plus a combinatorial
transport argument).
\end{enumerate}

Conversely, Fortini-style combinatorial transport lemmas are a template for proving ``hard'' equalities
of $\sigma$-algebras or probability expressions: reduce to finite approximants, prove invariance at
finite horizon, then pass to a limit (e.g.\ via monotone convergence).

\section{Recommended proof DAG (lemma-by-lemma) for a clean development}

A robust module DAG typically looks like:

\begin{enumerate}[label=\textbf{L\arabic*.}]
\item \textbf{(Base)} Definitions of cylinders, row processes, Markov parameters, and
generator-based measurable spaces on parameter types.
\item \textbf{(Measurability)} Lemmas proving measurability/AE-measurability of:
\begin{itemize}
\item row-process maps,
\item singleton evaluation maps of extracted kernels,
\item word probability compositions.
\end{itemize}
\item \textbf{(Row de Finetti)} Extraction of row kernels from per-row exchangeability.
\item \textbf{(Reconstruction infrastructure)} Change-of-variables (lintegral/map) lemmas and base cases
(nil/singleton).
\item \textbf{(Crux)} Cross-anchor product identity (or the minimal joint-law hypothesis that implies it).
\item \textbf{(Export)} The final representation theorem(s), plus clearly labeled variants:
\begin{itemize}
\item unrestricted Kleisli statement (if measurability is fully discharged),
\item Markov-only or ``safe endpoint'' statement (if not).
\end{itemize}
\end{enumerate}

\section{Positive and negative proof patterns}

\subsection{Positive pattern: explicit kernel mediator with factorization}
\begin{enumerate}[label=(\roman*)]
\item Define the mediator as a measurable map $f : Y \to \ProbMeas(\Theta)$, not merely pointwise.
\item Prove measurability by showing all coordinate maps $y \mapsto f(y)(A)$ are measurable for a
countable generating family of $A$.
\item Prove factorization by rewriting both sides as integrals of the same measurable integrand and
using standard measure-theory lemmas (e.g.\ $\int \mathbf{1}_E\,dP = P(E)$, monotone convergence,
$\mathrm{lintegral\_map}$).
\end{enumerate}

\subsection{Negative pattern: pointwise choice used as if measurable}
\begin{enumerate}[label=(\roman*)]
\item Use classical choice to pick $\pi_y$ from an existence statement.
\item Proceed to integrate against $y\mapsto \pi_y$ without proving measurability into $\ProbMeas(\Theta)$.
\item This compiles only if one silently changes the measurable structure, adds an axiom, or smuggles a
``measurable by simp'' lemma that is false.
\end{enumerate}

\section{Open problems and next steps (conceptual checklist)}

\begin{enumerate}[label=\arabic*.]
\item \textbf{Measurability upgrade lemma(s):} prove the missing inclusion between Borel and evaluation
$\sigma$-algebras under standard assumptions (often Polish/standard Borel).
\item \textbf{Fortini crux:} prove the cross-anchor product identity from existing invariances; if not
derivable, identify the minimal additional joint-law hypothesis and state it explicitly.
\item \textbf{Connection to universal prediction:} interpret the latent object (directing measure or
Markov parameter law) as a \emph{predictive state} that supports Bayesian updating by conditioning.
\item \textbf{PLN interface:} identify which objects are treated as ``belief states'' (measures on
hypotheses) and show how mixture representations support rule-based inference as measurable transformations.
\item \textbf{API hygiene:} keep a single canonical theorem per major endpoint, and re-export lightweight
aliases only when they remove real user friction.
\end{enumerate}

\section*{Acknowledgments}
This note is intentionally anonymous and conceptual; it summarizes patterns and bottlenecks that arise
when formalizing probabilistic representation theorems in proof assistants.

\end{document}
